\documentclass[10pt,a4paper]{amsart}
\usepackage[margin=19mm]{geometry}
\usepackage{amsmath,amssymb,mathtools}
\usepackage[scr=rsfs,cal=euler]{mathalfa}
\usepackage{xfrac}
\usepackage[unicode,hidelinks,bookmarks=false]{hyperref}
\newcommand{\ie}{i.e.\ }
\newcommand{\cf}{cf.\ }
\newcommand{\resp}{respectively\ }
\newcommand{\Subsection}{\subsection}
\providecommand{\nobreakdash}{\nobreak\hbox{-}}
\setlength{\emergencystretch}{2em}
\multlinegap0pt
\usepackage{xcolor}
\usepackage{mathtools}
\usepackage{bm}
\usepackage[shortlabels]{enumitem}
\newtheorem{proposition}{Proposition}
\newcommand{\X}{\mathcal{S}}
\title{A Finite-Iteration Theory for Asynchronous Categorical Distributional Temporal-Difference Learning}
\author{Ege C. Kaya, Abolfazl Hashemi}
\date{Source: arXiv:2605.06866v2 (CC BY 4.0)}
\begin{document}
\maketitle

\noindent\emph{Source and licence.} A Finite-Iteration Theory for Asynchronous Categorical Distributional Temporal-Difference Learning. Version arXiv:2605.06866v2 (CC BY 4.0), distributed by arXiv under the Creative Commons Attribution 4.0 International licence, http://creativecommons.org/licenses/by/4.0/, as recorded in the arXiv licence metadata for this exact version and shown on the abstract page https://arxiv.org/abs/2605.06866v2. Full licence notice: \texttt{licenses/arxiv-2605.06866-CC-BY-4.0.txt}.\par\smallskip

\noindent\textit{Editorial scope.} Counted unit: the article's proposition prop:mtd-proj-affine (source0.tex lines 1349--1372). The article's complete original proof of it is delivered with it (source0.tex lines 1374--1449, one \texttt{proof} environment of 76 lines). No mathematics is written, repaired or re-derived: the statement and the proof are the article's own lines, in the article's own notation, and no retained line is substituted or re-flowed.  No further source-local context is needed: every cross-reference of every retained line resolves inside the two delivered spans.  Nothing was omitted from any retained span, so the package carries no omission record.  No retained line contains a citation, so the wrapper carries no bibliography and no bibliography entry.  No retained line contains a figure environment, an image-inclusion command or drawing code, so no image is delivered and the package declares no asset dependency.  Editorial formatting only, in addition: an AMS \texttt{amsart} preamble that restates the article's own declarations and macros used by the retained lines, namely the environments \texttt{proposition} on separate counters, together with the standard packages those lines need.  The retained environments print in this document as Proposition 1 in order of appearance, whereas the article prints the same items with the numbering of its own full document.  The complete native source of the article as distributed in the arXiv e-print for this exact version is bundled unchanged as \texttt{sources/200056/source0.tex}.
\begin{proposition}[Embedded MTD projection as an affine Euclidean projector]
\label{prop:mtd-proj-affine}
For each state $s \in \X$, let
\begin{equation}
\mathcal A_s := K_s^{1/2}\mathbb R_1^d \subset \mathbb R^d .
\end{equation}
For every mass-$1$ finite signed measure $\nu$ on $\mathbb R^q$, define
\begin{equation}
b_s(\nu) := \left( \int \kappa(\theta_i(s),y)\,d\nu(y) \right)_{i=1}^d \in \mathbb R^d
\end{equation}
and
\begin{equation}
\zeta_s(\nu) := K_s^{\dagger/2} b_s(\nu),
\end{equation}
where $K_s^\dagger$ is the Moore-Penrose pseudoinverse of $K_s$. Then $b_s(\nu) \in \operatorname{range}(K_s)$ and
\begin{equation}
I_{\mathrm M,s}\bigl(\Pi_{\mathrm M}^{\Theta(s)}\nu\bigr) = \operatorname{proj}_{\mathcal A_s}\bigl(\zeta_s(\nu)\bigr),
\end{equation}
where $\operatorname{proj}_{\mathcal A_s}$ is the Euclidean orthogonal projection onto $\mathcal A_s$. Consequently, the map
\begin{equation}
J_s(\nu) := I_{\mathrm M,s}\bigl(\Pi_{\mathrm M}^{\Theta(s)}\nu\bigr)
\end{equation}
is affine on the affine space of mass-$1$ finite signed measures.
\end{proposition}

\begin{proof}
Let $\mathcal H_\kappa$ be the RKHS of $\kappa$, let $\varphi : \mathbb R^q \to \mathcal H_\kappa$ denote the feature map, and define the linear operator
\begin{equation}
\Phi_s : \mathbb R^d \to \mathcal H_\kappa, \qquad \Phi_s p := \sum_{i=1}^d p_i \varphi(\theta_i(s)).
\end{equation}
Then
\begin{equation}
K_s = \Phi_s^*\Phi_s .
\end{equation}
For a mass-$1$ finite signed measure $\nu$, define its kernel mean embedding by
\begin{equation}
m_\nu := \int \varphi(y)\,d\nu(y) \in \mathcal H_\kappa .
\end{equation}
Let $P_s$ denote the orthogonal projection in $\mathcal H_\kappa$ onto $\operatorname{range}(\Phi_s)$. Since $\Phi_s^*(I-P_s)=0$, we have
\begin{equation}
b_s(\nu) = \Phi_s^* m_\nu = \Phi_s^* P_s m_\nu \in \operatorname{range}(\Phi_s^*\Phi_s) = \operatorname{range}(K_s).
\end{equation}

Now fix $p \in \mathbb R_1^d$ and write
\begin{equation}
\mu_p := \sum_{i=1}^d p_i \delta_{\theta_i(s)} .
\end{equation}
By the reproducing-kernel representation of MMD,
\begin{equation}
\mathrm{MMD}_\kappa(\mu_p,\nu)^2 = \lVert \Phi_s p - m_\nu \rVert_{\mathcal H_\kappa}^2 = \lVert \Phi_s p - P_s m_\nu \rVert_{\mathcal H_\kappa}^2 + \lVert (I-P_s)m_\nu \rVert_{\mathcal H_\kappa}^2 .
\end{equation}

We next identify $P_s m_\nu$ in coefficient form. Since $K_s=\Phi_s^*\Phi_s$, the standard Moore-Penrose identity gives
\begin{equation}
\Phi_s K_s^\dagger \Phi_s^* = P_s .
\end{equation}
Therefore,
\begin{equation}
\Phi_s\bigl(K_s^\dagger b_s(\nu)\bigr) = \Phi_s K_s^\dagger \Phi_s^* m_\nu = P_s m_\nu .
\end{equation}
Substituting this into the previous display yields
\begin{equation}
\lVert \Phi_s p - P_s m_\nu \rVert_{\mathcal H_\kappa} = \lVert \Phi_s(p-K_s^\dagger b_s(\nu)) \rVert_{\mathcal H_\kappa}.
\end{equation}
Since $\lVert \Phi_s x \rVert_{\mathcal H_\kappa}^2 = x^\top K_s x = \lVert K_s^{1/2}x \rVert_2^2$ for every $x \in \mathbb R^d$, we obtain
\begin{equation}
\lVert \Phi_s p - P_s m_\nu \rVert_{\mathcal H_\kappa} = \lVert K_s^{1/2}(p-K_s^\dagger b_s(\nu)) \rVert_2 = \lVert K_s^{1/2}p - K_s^{1/2}K_s^\dagger b_s(\nu) \rVert_2 .
\end{equation}
Because $K_s$ is symmetric positive semidefinite, spectral calculus gives
\begin{equation}
K_s^{1/2}K_s^\dagger = K_s^{\dagger/2}.
\end{equation}
Hence
\begin{equation}
\lVert \Phi_s p - P_s m_\nu \rVert_{\mathcal H_\kappa} = \lVert K_s^{1/2}p - K_s^{\dagger/2} b_s(\nu) \rVert_2 = \lVert K_s^{1/2}p - \zeta_s(\nu) \rVert_2 .
\end{equation}
Therefore
\begin{equation}
\mathrm{MMD}_\kappa(\mu_p,\nu)^2 = \lVert K_s^{1/2}p - \zeta_s(\nu) \rVert_2^2 + c_s(\nu),
\end{equation}
where
\begin{equation}
c_s(\nu) := \lVert (I-P_s)m_\nu \rVert_{\mathcal H_\kappa}^2
\end{equation}
does not depend on $p$.

Minimizing over $p \in \mathbb R_1^d$ is therefore equivalent to Euclidean projection of $\zeta_s(\nu)$ onto
\begin{equation}
\mathcal A_s = K_s^{1/2}\mathbb R_1^d .
\end{equation}
Since $I_{\mathrm M,s}\bigl(\sum_{i=1}^d p_i\delta_{\theta_i(s)}\bigr)=K_s^{1/2}p$, it follows that
\begin{equation}
I_{\mathrm M,s}\bigl(\Pi_{\mathrm M}^{\Theta(s)}\nu\bigr) = \operatorname{proj}_{\mathcal A_s}\bigl(\zeta_s(\nu)\bigr).
\end{equation}

Finally, $\nu \mapsto b_s(\nu)$ is linear, hence so is $\nu \mapsto \zeta_s(\nu)=K_s^{\dagger/2}b_s(\nu)$. Since orthogonal projection onto an affine subspace is an affine map, the composition
\begin{equation}
J_s(\nu) = \operatorname{proj}_{\mathcal A_s}\bigl(\zeta_s(\nu)\bigr)
\end{equation}
is affine on the affine space of mass-$1$ finite signed measures.
\end{proof}

\end{document}
